\documentclass[11pt]{article}
\usepackage{amsmath,amssymb,hyperref}
\title{What a certified retirement ledger forgets about formerly forbidden orders}
\date{25 November 2026}
\begin{document}
\maketitle
\paragraph{Retained law.}
Use the physical effect-side retirement and fixed signature rows of
\href{https://github.com/eternal-november/eos-research-programmes/pull/384}{PR 384}.
Now a scope is lawful only after its prerequisites have actually settled. Assume these prerequisites form a fixed finite partial order \(P\), each scope runs once, and every linear extension is physically and lawfully available in the test language. A claim carries a fixed footprint and a fixed set \(A\) of scopes, any one of which retires it when actually certified. Capacity tests see the accumulated remaining footprint at a settled cut. No authorization, request or promise is treated as a certified effect. Each conclusion applies coordinatewise to the retained footprint vector. Time, further changes to signatures, extra receipt fields and state-dependent reachability are outside the model.
\paragraph{Exact old profile.}
The completed sets are exactly the lower ideals \(I\) of \(P\); write \(U=P\setminus I\) for their complementary upsets. A claim survives iff \(A\subseteq U\), equivalently \(\uparrow_P A\subseteq U\). Thus sum the raw weights by \(B=\uparrow_P A\) to obtain \(w_B\). The answers on all lawful tests are
\[
      f(U)=\sum_{B\subseteq U}w_B.
\]
Every upset in this sum can be the complement of an actual settled prefix in the declared test language. Recursively subtract terms for strict subsets to recover \(w_U\) from those answers. This is the ordinary
\href{https://webhomes.maths.ed.ac.uk/~v1ranick/papers/rota1.pdf}{finite incidence transform}.
The grouped profile therefore characterizes equality of \emph{all} lawful continuation responses, not only the elementary retirements; it also suffices for adaptively selected continuations, since each successive settled cut is still such an ideal.
\paragraph{An amendment that opens a new test.}
Let \(Q\) have the same scopes and weaker prerequisites, leaving actual release rules and raw rows unchanged. The new bins map onto old ones by \(B\mapsto\uparrow_P B\), since \(\uparrow_P(\uparrow_Q A)=\uparrow_P A\). Strengthening alone can sum bins along this map; weakening must split them. Require \(a\) before \(b\). A one-unit claim tagged \(\{a\}\) and one tagged \(\{a,b\}\) have identical old profiles. If an authorized amendment permits \(b\) first, the second claim releases and the first does not. An actual new job needing the unit can run only after the former effect. Checking \(b\)'s permission without checking its completed effect still does not establish capacity.
\par
More generally, choose one newly admitted upset \(U_*\). Assign its answer one and every other new-law answer zero, then invert to get signed integer weights. Splitting these into positive and negative ledgers gives two nonnegative histories that agree on every old response and on every new response except \(U_*\). Their shortest distinguishing operation count is precisely \(|Q\setminus U_*|\). In a separate three-scope example with \(b<c\), the two ledgers \(\{\varnothing,\{a,c\}\}\) and \(\{\{a\},\{c\}\}\), one unit each, agree on every proper reachable ideal and differ only after all three effects.
\paragraph{How much can a single prerequisite hide?}
Suppose \(P\) is obtained from \(Q\) by adding one ordered pair \(a<b\) (and its transitive consequences), with \(a,b\) incomparable in \(Q\). A newly lawful upset must contain \(a\) but not \(b\). It therefore contains all of \(\uparrow_Q a\) and omits all of \(\downarrow_Q b\); these sets are disjoint. Its unconstrained remainder is any upset of the induced subposet
\[
      R=Q\setminus (\uparrow_Q a\cup\downarrow_Q b).
\]
Conversely each such remainder gives an upset of \(Q\): a required upward step cannot enter the excluded downset from outside it. Thus the number \(D\) of new scalar answer coordinates is exactly the number of upsets of \(R\), at most \(2^{n-2}\) for \(n\) scopes. An antichain \(Q\) attains the bound. This counts the whole new interface, not the information needed for one named downstream job.
\par
For arbitrary real nonnegative weights, one scalar footprint coordinate and the old full profile given, \(D\) additional independent scalar capacity tests are necessary and sufficient to reconstruct the new full profile in the resettable query model. Sufficiency is to query all the new upsets and invert. For necessity use a strictly positive new profile, which is attainable by raw signatures equal to those upsets. After fewer than \(D\) tests even an adaptive policy leaves a nonzero perturbation in the kernel of the old and tested rows, by the full-rank transform. A small perturbation in that direction stays positive and returns exactly the same answers on the chosen query path, but has different new profile. This lower bound concerns scalar tests, not bits or unlimited encodings. A real running ledger cannot in general reset and execute incompatible paths; observing one live sequence of settlements gives at most \(n\) further cut values. In particular when \(D>n\), even a complete live sequence cannot in general reconstruct the whole new profile under these observation rules. Store more beforehand or certify just the queries that the new job will consume.
\paragraph{Boundary and checks.}
The direct ideal enumeration agrees. The replay independently constructs raw rows, projects the new profile to the old one and recovers test profiles by inversion, checks the exact edge-change count, and gives a separate signed-inversion witness for each newly lawful test. Two Lean lemmas check closure transport and survival. The native ordered prefix in \href{https://github.com/eternal-november/eos-research-programmes/pull/441}{PR 441} supports separating actual acceptance from carrier availability; it does not supply a real authority to remove a prerequisite. The amendment and numerical ledger are an authored proof boundary, not a native replay, independent transfer, or evidence of faster service.
\end{document}
